% TOP_PROBLEM: {"id": 3214, "title": "5-flow conjecture", "category_id": 3, "category": {"id": 3, "name": "graph_theory", "display_name": "Graph Theory", "description": "Problems involving graphs, networks, and their properties.", "slug": "graph-theory", "order_index": 3, "created_at": "2026-07-31T15:26:25.671Z"}, "status": "open", "problem_number": "OPG-126", "set_id": 12}
% FIRST_POSTED: 2026-10-01
% ATTEMPT_STATUS: unresolved
% MODEL: GPT 6 Astra Ultra
% UnsolvedMath ID 3214; source catalogue code OPG-126
% !TeX program = lualatex
% Research attempt dated 2026-10-01; canonical statement preserved.
\documentclass[11pt,a4paper]{article}
\usepackage[margin=25mm]{geometry}
\usepackage{fontspec}
\setmainfont{lmroman10-regular.otf}[BoldFont=lmroman10-bold.otf,
 ItalicFont=lmroman10-italic.otf,BoldItalicFont=lmroman10-bolditalic.otf]
\usepackage{amsmath,amssymb,mathtools}
\usepackage{unicode-math}
\setmathfont{latinmodern-math.otf}
\AtBeginDocument{\let\setminus\smallsetminus}
\usepackage{xurl,hyperref}
\hypersetup{colorlinks=true,urlcolor=blue}
\setcounter{secnumdepth}{0}
\setlength{\parindent}{0pt}
\setlength{\parskip}{5pt}
\setlength{\emergencystretch}{3em}
\begin{document}
\section*{5-flow conjecture}\label{3214}
{\small UnsolvedMath ID 3214 · OPG-126 · Graph Theory · OpenGarden}
\subsection{Definitions and mathematical statement}
Let $G=(V,E)$ be a finite undirected
loopless multigraph; parallel edges
are distinguished. A bridge is an
edge whose deletion increases the
number of connected components.
Choose an orientation of every edge.
Write $\delta^+(v)$ and $\delta^-(v)$
for the edges directed out of and
into vertex $v$, respectively.

For an integer $q\geq2$, an integer
$q$-flow is a function $\phi:E\to\mathbb Z$
satisfying conservation at every vertex,
\[
 \sum_{e\in\delta^+(v)}\phi(e)
 =\sum_{e\in\delta^-(v)}\phi(e),
 \qquad |\phi(e)|<q\quad(e\in E).
\]
It is nowhere-zero if additionally
$\phi(e)\neq0$ on every edge.
Reversing the orientation of one
edge and negating its value preserves
these conditions, so existence is
independent of the initially chosen
orientation.

Tutte's $5$-flow conjecture asserts
that every bridgeless $G$ admits a
nowhere-zero integer $5$-flow. Thus
each edge must carry one of
$\pm1,\pm2,\pm3,\pm4$, with exact
integer conservation at every vertex.
Equivalently, an orientation may be
chosen so all values belong to
$\{1,2,3,4\}$.

Bridgelessness is necessary: summing
conservation over one component after
deleting a bridge forces its value
to be zero. Components and isolated
vertices cause no additional condition;
they may be handled independently.
The bound five is uniform over all
orders and degrees, rather than a
bound allowed to depend on the graph.
\subsection{Short English statement}
Can every bridgeless graph carry a conserved nonzero integer flow with each edge value having magnitude at most four?
\subsection{Research attempt}
\paragraph{Scope and method.}
This research attempt by GPT-6 Astra Ultra is dated 1 October 2026.
It preserves the catalogue statement and distinguishes elementary proofs
below from cited prior work. No novelty claim or formal proof-assistant
verification is made. The final paragraph states the precise remaining scope.
\paragraph{Literature and the boundedness issue.}
The primary paper [R1] gives proofs of Seymour's theorem that every
bridgeless graph has a nowhere-zero six-flow. That is a prior theorem,
not a five-flow proof. We first give an elementary construction showing
why bridges are the exact obstruction to an integer flow with some finite
bound, then obtain the required bound on two substantial special classes.
All conservation equations below are integer equations, not merely
congruences modulo five.

\paragraph{Bridges versus cycle coverage.}
If $e$ is a bridge, sum conservation over one component of $G-e$.
All internal contributions cancel, leaving $\pm\phi(e)=0$, so a nowhere-zero
flow is impossible. Conversely, every edge in a bridgeless loopless
multigraph lies on a cycle: after deleting it there remains a path between
its endpoints. A pair of parallel edges can form a cycle of length two.
Choose finitely many cycles $C_0,\ldots,C_{r-1}$ covering all edges and
assume here that $E(G)\ne\varnothing$, so $r\ge1$; an edgeless graph
already has the empty two-flow. Orient each chosen cycle cyclically.
Relative to one fixed reference orientation, its
unit circulation is a vector $z_i\in\{-1,0,1\}^{E}$.
Put
\[
             \phi=\sum_{i=0}^{r-1}2^i z_i.
\]
Each vector satisfies conservation, so their sum does. For an edge $e$,
let $j$ be the largest index with $z_j(e)\ne0$. Its contribution has
magnitude $2^j$, strictly greater than
$\sum_{i<j}2^i=2^j-1$. Hence $\phi(e)\ne0$, and
$|\phi(e)|\le2^r-1$. We have constructed a $2^r$-flow without using a
general bounded-flow theorem. The exponential bound explains exactly what
this construction fails to control.

\paragraph{Eulerian components.}
If all degrees are even, every component with edges has a closed Euler
tour. Orient its edges along that tour and assign value one. Every arrival
is paired with a departure, so conservation holds at every vertex.
Equivalently, decompose the edges into cycles and orient them separately;
edge-disjointness avoids cancellation altogether. This is a nowhere-zero
two-flow and therefore satisfies the five-flow requirement. Isolated
vertices impose the empty equation and need no separate adjustment.

\paragraph{Three-edge-colourable cubic graphs.}
Let the colour classes of a proper three-edge-colouring be $M_1,M_2,M_3$.
Every vertex meets one edge of each colour. The graphs with edge sets
$M_1\cup M_2$ and $M_2\cup M_3$ are disjoint unions of cycles, so admit
unit circulations $x,y$, extended by zero outside their respective supports.
The integer circulation
\[
                         f=x+2y
\]
has magnitude one on $M_1$, magnitude two on $M_3$, and magnitude one
or three on $M_2$: there its value is $\pm1\pm2$. Thus it never vanishes
and has magnitude at most three. This proves a four-flow, stronger than
the requested five-flow, for every such cubic graph. The argument works
for parallel edges as well because the two-colour components are still
even cycles, including length two.

\paragraph{Why the construction does not extend automatically.}
A bridgeless cubic graph need not admit three edge colours. An arbitrary
cycle cover can have many overlapping supports, and assigning all cycles
weight one can produce cancellations. Powers of two prevent cancellation
but lose a uniform bound. Even finding two spanning even subgraphs is
not automatic: their union must cover every edge for the displayed
$x+2y$ argument to be nowhere zero. The proof does not establish that
coverage property for an arbitrary bridgeless graph.

\paragraph{Verification and final status.}
The root batch script verifies a concrete $K_4$ certificate by computing
the signed incidence sums of both circulations and their weighted sum.
The bridge obstruction and largest-term estimate are exact symbolic
arguments valid at every size. \textbf{UNRESOLVED.} Eulerian graphs and
three-edge-colourable cubic graphs satisfy the target here. Controlling
overlapping cycle circulations by the fixed magnitudes one through four
in every bridgeless graph is the missing universal step.

\subsection{Sources}
\begin{itemize}
\item[S1] ULAM.AI and the UnsolvedMath Contributors, supplied problems.json, record 3214 (OPG-126), OpenGarden collection. Catalogue identity and source transcription; CC BY 4.0.
\url{https://huggingface.co/datasets/ulamai/UnsolvedMath}.
\item[S2] Open Problem Garden, 5-flow conjecture. Original problem and discussion; the mathematical formulation below is expanded from the supplied source record.
\url{http://www.openproblemgarden.org/op/5_flow_conjecture}.
\item[R1] Matt DeVos, Edita Rollová and Robert Šámal,
\emph{A new proof of Seymour's 6-flow theorem}, arXiv:1512.06214.
Primary abstract checked 1 October 2026.
\url{https://arxiv.org/abs/1512.06214}.
\end{itemize}
\end{document}
